\clearpage\begingroup\SourceLeaf{58b}\HistoricalHead{53}{R}{GAUSS AN ENCKE.}
\begin{center}
[5.]\\[.7em]
\textsc{Gauss} an \textsc{Schumacher}[\textsuperscript{**})].
\par\smallskip\rule{.78\textwidth}{.35pt}\\
{\small [Briefwechsel zwischen \textsc{Gauss} und \textsc{Schumacher} II, Altona 1860, S. 220.]}\\
\rule{.78\textwidth}{.35pt}
\end{center}
\begingroup\renewcommand{\thefootnote}{}\footnotetext{[\textsuperscript{**}) Das Datum fehlt, muß aber zwischen dem 7. und dem 30. Dezember 1829 liegen.]}\endgroup
\noindent .\ .\ .\ .\ .\ .\ .\ .\ .\ .\ .\ .\ .\ .\ .\ .\ .\ .\ .\ .\ .
\par\noindent\hspace*{\parindent}Vor einiger Zeit zeigten Sie in den A[stronomischen] N[achrichten][\textsuperscript{***})]\begingroup\renewcommand{\thefootnote}{}\footnotetext{[\textsuperscript{***}) Astronom. Nachrichten \textbf{7}, 1829, Spalte 103.]}\endgroup\ an, dass Sie aus England eine Linse erhalten hätten, um die \textsc{Brown}schen
Versuche über das Leben der Materie zu wiederholen. Was ist wol das Resultat
dieser Wiederholung gewesen.\quad .\ .\ .\ .\ .\ .\ .\ .\ .\ .\ .\ .\ .\ .\ .\ .\ .\ .\ .\ .\ .
\clearpage\endgroup
